\documentclass[10pt]{article}
\usepackage[T1]{fontenc}
\usepackage[utf8]{inputenc}
\usepackage[french]{babel}
\usepackage{amsmath,amssymb,mathrsfs,mathpazo}
\usepackage{geometry}
\geometry{a4paper,margin=2cm}
\usepackage{graphicx,array,tabularx,tabularray,colortbl,multirow,multicol,booktabs,enumitem,pifont,ulem,fancybox,diagbox,cancel}
\usepackage{pstricks,pstricks-add,pst-plot,pst-tree,pst-node,pst-text,pst-func,pst-math,pst-bspline,pst-3dplot,pst-eucl,pst-solides3d}
\let\pscalculate\undefined
\usepackage{pst-bezier}
\usepackage{tikz,pgfplots,tkz-tab}
\usetikzlibrary{arrows,arrows.meta,patterns,calc,positioning,trees,decorations.markings,shapes.geometric,shapes.misc,shapes.arrows,intersections,decorations.pathmorphing,decorations.pathreplacing,backgrounds,shadows,babel}
\pgfplotsset{compat=1.18}
\usepgfplotslibrary{fillbetween}
\usepackage[np]{numprint}
\usepackage{listings,makecell,pas-tableur,twoopt}
\let\barre\undefined
\usepackage[pstricks]{tablvar}
\definecolor{darkgreen}{RGB}{0,100,0}
\providecommand{\red}{\color{red}}
\providecommand{\blue}{\color{blue}}
\providecommand{\green}{\color{green}}
\providecommand{\black}{\color{black}}
\providecommand{\white}{\color{white}}
\providecommand{\gray}{\color{gray}}
\providecommand{\cyan}{\color{cyan}}
\providecommand{\magenta}{\color{magenta}}
\providecommand{\yellow}{\color{yellow}}
\providecommand{\violet}{\color{violet}}
\providecommand{\orange}{\color{orange}}
\providecommand{\euro}{\texteuro}
\providecommand{\EUR}{\texteuro}
\providecommand{\ieme}{\textsuperscript{e}}
\providecommand{\degres}{\ensuremath{^\circ}}
\providecommand{\degree}{\ensuremath{^\circ}}
\providecommand{\vv}[1]{\overrightarrow{#1}}
\providecommand{\Coord}[1]{#1}
\providecommand{\decosix}{}
\providecommand{\starredbullet}{\ensuremath{\bullet}}
\pagestyle{empty}
\setlength{\parindent}{0pt}
\usepackage{adjustbox}\DeclareUnicodeCharacter{2236}{:}\DeclareUnicodeCharacter{2190}{\ensuremath{\leftarrow}}
\providecommand{\R}{\mathbb{R}}
\providecommand{\N}{\mathbb{N}}
\providecommand{\Z}{\mathbb{Z}}
\providecommand{\Q}{\mathbb{Q}}
\providecommand{\D}{\mathbb{D}}
\providecommand{\C}{\mathbb{C}}
\providecommand{\e}{{\mathrm{e}}}
\providecommand{\vect}[1]{\overrightarrow{#1}}
\providecommand{\vectt}[1]{\overrightarrow{\mathrm{#1}}}
\providecommand{\vvt}[1]{\overrightarrow{\mathrm{#1}}}
\providecommand{\np}[2][]{{#2~#1}}
\providecommand{\vv}[1]{{\overrightarrow{#1}}}
\renewcommand{\d}{{\mathrm{d}}}
\providecommand{\degres}{{^\circ}}
\providecommand{\degre}{{^\circ}}
\providecommand{\blue}{\color{blue}}
\providecommand{\red}{\color{red}}
\providecommand{\enskip}{\quad}
\providecommand{\llbracket}{\lbrack\!\lbrack}
\providecommand{\rrbracket}{\rbrack\!\rbrack}
\providecommand{\euro}{{\text{€}}}
\providecommand{\no}{{\text{n°}}}
\providecommand{\interval}[2]{{[#1;#2]}}
\providecommand{\intervFO}[2]{{[#1;#2[}}
\providecommand{\dotfill}{\ldots}

\providecommand{\ds}{\displaystyle}
\providecommand{\textding}[1]{#1}
\providecommand{\Oij}{\(\mathrm{O};\vec{\imath},\vec{\jmath}\)}
\providecommand{\Oijk}{\(\mathrm{O};\vec{\imath},\vec{\jmath},\vec{k}\)}
\providecommand{\Ouv}{\(\mathrm{O};\vec{u},\vec{v}\)}
\providecommand{\intervFF}[2]{\ensuremath{[#1;#2]}}
\renewcommand{\euro}{\ifmmode\text{\texteuro}\else\texteuro\fi}
\DeclareUnicodeCharacter{00B0}{\ifmmode^{\circ}\else\textdegree\fi}
\DeclareUnicodeCharacter{20AC}{\ifmmode\text{\texteuro}\else\texteuro\fi}
% Énoncés : Éducation nationale. Transcription et corrigés : APMEP, Denis Vergès et contributeurs. Découpage et classement : CoopMaths, Sébastien Lozano, Éric Elter et contributeurs. Adaptation de la présentation : Maths Explorer.
% CC BY-SA : https://coopmaths.fr/www/about/
% https://coopmaths.fr/alea/static/e3c/2020/tex/e3c_2020_00_sujet52_3_cor.tex
\begin{document}
\noindent\textit{APMEP / CoopMaths. CC BY-SA. Présentation adaptée par Maths Explorer.}\par\medskip
\(f(x) = 60x \e^{-0{,}5x}\).

\subsection*{1.}

\( f \) est un produit de fonctions dérivables sur \( \mathbb{R} \), donc sur \([0 \,;\, 10]\) et sur cet intervalle :
\begin{align*}
f'(x) &= 60\e^{-0{,}5x} + 60x \times (-0{,}5)\e^{-0{,}5x} \\
&= \e^{-0{,}5x}(60 - 30x) \\
&= -30(x - 2)\e^{-0{,}5x}.
\end{align*}

\subsection*{2.}

On sait que, quel que soit le réel \( x \), \( \e^{-0{,}5x} > 0 \), donc le signe de \( f'(x) \) est celui de \( 2 - x \) :
\begin{align*}
&2 - x > 0 \\
\iff &2 > x \\
\iff &x < 2.
\end{align*}
Sur \([0 \,;\, 2]\), \( f'(x) > 0 \), sur \([2 \,;\, 10]\), \( f'(x) < 0 \) et \( f'(2) = 0 \).

\subsection*{3.}

D'après la question précédente \( f \) est croissante sur \([0 \,;\, 2]\) de \( f(0) = 0 \) à \( f(2) = 120\e^{-1} \approx 44{,}1 \), puis décroissante sur \([2 \,;\, 10]\) de \( f(2) = 120\e^{-1} \) à \( f(10) = 600\e^{-5} \approx 4{,}04 \). D'où le tableau :

\begin{center}
\begin{tikzpicture}
\tkzTabInit[lgt=3.5, espcl=4]{$x$ / 1, {Signe de $f'(x)$} / 1, {$f(x)$} / 2}{${0}$, ${2}$, ${10}$}
\tkzTabLine{,+,0,-,}
\tkzTabVar{-/{$0$},+/{$120\e^{-1}$},-/{$600\e^{-5}$}}{/}
\end{tikzpicture}
\end{center}

\subsection*{4.}
On a vu que \( f'(2) = 0 \) : le coefficient directeur de la tangente à la courbe représentative de la fonction \( f \) est parallèle à l'axe des abscisses (horizontale).

\subsection*{5.}

On sait qu'une équation réduite de la tangente à la courbe représentative de la fonction \( f \) au point d'abscisse 0 est :
\[
y - f(0) = f'(0)(x - 1).
\]
Avec :
\[
f(0) = 0 \quad \text{et} \quad f'(0) = -30(0 - 2)\e^{-0{,}5 \times 0} = 60 \e^0 = 60 \times 1 = 60,
\]
l'équation devient \(y = 60x\).

\end{document}
